\documentclass[12pt]{article}
\usepackage[brazil]{babel}
\usepackage[utf8]{inputenc}
\usepackage[T1]{fontenc}
\usepackage{amsmath, amssymb, amsthm}
\usepackage{geometry}
\usepackage{graphicx}
\usepackage{xcolor}
\usepackage{enumitem}
\usepackage{listings}

\definecolor{maincolor}{RGB}{37, 150, 190}
\usepackage{fancyhdr} 
\definecolor{forestgreen}{RGB}{37, 150, 190}
\usepackage[most]{tcolorbox}

\lstset{
  basicstyle=\ttfamily\small,
  keywordstyle=\color{blue},
  stringstyle=\color{red},
  commentstyle=\color{green!50!black},
  breaklines=true,
  frame=single,
  numbers=left,
  numberstyle=\tiny\color{gray},
}

\newtcolorbox{block}[1]{%
  colback=green!5,
  colframe=maincolor,
  coltitle=white, 
  title=#1,
  fonttitle=\bfseries,
  boxrule=1.5pt,
  arc=3mm,
  left=6pt,
  right=6pt,
  top=6pt,
  bottom=6pt
}

\geometry{margin=2.5cm}
\pagestyle{fancy}
\fancyhf{}

\renewcommand{\headrulewidth}{2pt}
\renewcommand{\footrulewidth}{2pt}

\renewcommand{\headrule}{\hbox to\headwidth{%
  \color{forestgreen}\leaders\hrule height \headrulewidth\hfill}}

\renewcommand{\footrule}{\hbox to\headwidth{%
  \color{forestgreen}\leaders\hrule height \footrulewidth\hfill}}

\setlength{\headheight}{32pt}

\fancyhead[L]{%
  \includegraphics[height=1.2cm]{../../aulas/4.Relacoes/logo_novo.png}
}

\fancyhead[R]{%
  \raggedleft
  \textcolor{forestgreen}{%
    \small
    \textbf{Atividade de Programação.}\\
    Me.~Sergio Souza Novak\\
    \textit{Lógica para Engenharia de Software}
  }
}

\fancyfoot[C]{%
  \textcolor{forestgreen}{\thepage}
}

\begin{document}

\vspace*{0.2cm}
\section*{Exercícios de Programação — Relações e Matrizes}

\begin{block}{Instruções Gerais}
Para cada um dos exercícios abaixo, escolha uma das seguintes linguagens de programação: \textbf{C/C++ ou Python}. 
\textbf{Não serão aceitas atividades em JAVASCRIPT};

\textbf{Regras de Entrada:}
\begin{itemize}
    \item O programa deve ler um conjunto $A$ a partir do teclado (elementos únicos).
    \item Os elementos dos conjuntos devem ser \textbf{números naturais} ($0, 1, 2, \dots$).
    \item A leitura deve ser interrompida quando o usuário digitar um \textbf{número menor que zero} (ex: -1).
    \item A relação $R$ deve ser lida como pares de elementos do conjunto $A$.
\end{itemize}

\textbf{Formato de Entrega:}
\begin{itemize}
    \item A entrega deve ser um arquivo \textbf{.zip}.
    \item Dentro do arquivo .zip, deve haver \textbf{uma pasta para cada exercício} (Ex: Exercicio1, Exercicio2, etc.).
    \item Dentro de cada pasta, devem estar contidos \textbf{todos os códigos e arquivos necessários} para compilar e executar a resposta.
\end{itemize}
\end{block}

\vspace{0.4cm}
\begin{enumerate}

\item \textbf{Representação por Matriz de Relação}\\
Implemente um programa que leia o conjunto $A$ e uma relação $R \subseteq A \times A$. O programa deve gerar e exibir a \textbf{Matriz de Relação} $M_R$. A matriz deve ser quadrada ($n \times n$, onde $n$ é o número de elementos em $A$).

\begin{block}{Exemplo}
\textbf{Entrada:} \\
Conjunto A: 1, 2, 3, -1 \\
Relação R: (1, 1), (1, 2), (2, 3), -1 \\
\textbf{Saída:} \\
Matriz $M_R$: \\
$\begin{pmatrix} 1 & 1 & 0 \\ 0 & 0 & 1 \\ 0 & 0 & 0 \end{pmatrix}$
\end{block}

\vspace{0.3cm}
\item \textbf{Verificação de Reflexividade e Simetria}\\
Implemente um algoritmo que, dada a Matriz de Relação $M_R$ (ou a relação $R$ diretamente), determine se a relação é:
\begin{enumerate}[label=\alph*)]
    \item \textbf{Reflexiva:} $\forall a \in A, (a, a) \in R$.
    \item \textbf{Simétrica:} $\forall a, b \in A, (a, b) \in R \implies (b, a) \in R$.
\end{enumerate}

\begin{block}{Exemplo}
\textbf{Entrada:} \\
Conjunto A: 1, 2, -1 \\
Relação R: (1, 1), (2, 2), (1, 2), (2, 1), -1 \\
\textbf{Saída:} \\
Reflexiva: Sim \\
Simétrica: Sim
\end{block}

\vspace{0.3cm}
\item \textbf{Verificação de Transitividade}\\
Implemente um código que verifique se a relação $R$ sobre o conjunto $A$ é \textbf{transitiva}. Uma relação é transitiva se $\forall a, b, c \in A, ((a, b) \in R \land (b, c) \in R) \implies (a, c) \in R$.

\begin{block}{Exemplo}
\textbf{Entrada:} \\
Conjunto A: 1, 2, 3, -1 \\
Relação R: (1, 2), (2, 3), (1, 3), -1 \\
\textbf{Saída:} \\
Transitiva: Sim \\
\textbf{Entrada 2:} \\
Relação R: (1, 2), (2, 3), -1 \\
\textbf{Saída 2:} \\
Transitiva: Não (falta (1, 3))
\end{block}

\vspace{0.3cm}
\item \textbf{Relação de Equivalência}\\
Utilizando as verificações dos exercícios anteriores, implemente um programa completo que diga se uma relação $R$ sobre $A$ é uma \textbf{Relação de Equivalência}. Lembre-se que, para ser de equivalência, a relação deve ser simultaneamente Reflexiva, Simétrica e Transitiva.

\begin{block}{Exemplo}
\textbf{Entrada:} \\
Conjunto A: 1, 2, 3, -1 \\
Relação R: (1, 1), (2, 2), (3, 3), (1, 2), (2, 1), -1 \\
\textbf{Saída:} \\
Reflexiva: Sim \\
Simétrica: Sim \\
Transitiva: Sim \\
Relação de Equivalência: Sim
\end{block}

\end{enumerate}

\end{document}
